\chapter{Внутренний спектр планкона}
\label{ch:floor0-spectrum}

\section{Объект и этаж}
\label{sec:floor0-object}

Глава описывает \emph{этаж~0}: один планкон в одной планковской ячейке~$\hv=\PlanckCell$.
Занятость $b=1$, топологический заряд $n=\pm1$ на контуре (\cref{sec:planckon-floor0,ch:matter}).
Это не картина «орбитали внутри радиуса $r<\ell_P$'' --- ниже шага~$\hl$ пространства как континуума нет.
И не этаж~2 (кварки, составные квантовые числа): здесь только внутренние степени свободы
\emph{одного} локализованного возбуждения.

На макроуровне~$T$ атомные уровни и орбитали --- readout составного узла;
на~$\Mlayer$ у планкона заранее есть та же \emph{логика запрещёнки},
что у атома или ядра: дискретные уровни, правила отбора, запрет Паули,
конечный набор состояний (бит-бюджет $B_{hV}$, \cref{ch:evolution}).

\begin{table}[htbp]
  \centering
  \caption{Аналогия: атом / ядро на~$T$ и планкон на~$\Mlayer$ (этаж~0).}
  \label{tab:floor0-analogy}
  \begin{tabular}{@{}ll@{}}
    \toprule
    Атом, ядро ($T$) & Планкон ($\Mlayer$) \\
    \midrule
    уровни $E_n$ & $E=n_E E_0$, $n_E\in\mathbb{Z}$; фаза $\Phi$ на кольце \\
    орбитали $(n,\ell,m)$ & внутренние: кольцо + $SU(2)$; внешние: оболочки соседей (гл.~\ref{ch:floor0-spectrum}, \S\ref{sec:floor0-open}) \\
    запрет Паули & A16: два совпадающих спинора в одном $\hv$ запрещены; $b\in\{0,1\}$ \\
    запрещённый переход & скачок не кратен $E_0$, $\Delta\varphi_{\min}$ или ломает $n$ без партнёра \\
    запрещённая зона (кристалл) & щели $\omega(k)$ на зоне Бриллюэна вакуума (\cref{ch:matter}) \\
    \bottomrule
  \end{tabular}
\end{table}

\section{Два внутренних регистра}
\label{sec:floor0-registers}

На одной ячейке независимых «полей Стандартной модели» нет.
Из аксиом и закона~$g$ (\cref{ch:axioms,ch:evolution}) следуют два связанных регистра:

\paragraph{Фазовое кольцо.}
Накопленная фаза импульсного реестра
$\Phi\in\mathbb{Z}_{N_{\mathrm{ring}}}$, $N_{\mathrm{ring}}=512$.
Один полный оборот $U(1)$ на кольце --- $512$ дискретных тиков.
Минимальный разрешённый скачок $\Delta\varphi_{\min}=\tfrac12$ рад
(\cref{lem:51a}) даёт
\[
  \Delta\varphi_{\mathrm{disc}}
  = \left\lfloor N_{\mathrm{ring}}\,\frac{\Delta\varphi_{\min}}{2\pi}\right\rfloor
  = 41
\]
тик на $\mathbb{Z}_{512}$.
Один такой скачок --- один квант энергии $E_0$ лестницы (\cref{thm:51,lem:51e}).

\paragraph{Спинор на ядре.}
Состояние $z\in\mathbb{C}^2$ в центре~$\hv$.
Вращение вокруг оси Блоха ядра реализует $SU(2)$:
при повороте на $2\pi$ состояние меняет знак, на $4\pi$ --- возвращается
(полуцелый спин, A16).
Это \emph{внутренняя} ориентация, не топологический заряд $n$ на контуре.

Оба регистра живут на одном носителе и ограничены одним бит-бюджетом;
их нельзя трактовать как два произвольных квантовых числа лабораторной таблицы.

\section{Лестница $n_E$ на кольце}
\label{sec:floor0-ne-ladder}

Энергия события на~$\Mlayer$ квантуется лестницей $E=n_E E_0$ (\cref{lem:51e}).
Число $n_E$ считается по накопленным тикам фазы на кольце:
\[
  n_E = \left\lfloor \frac{|\Phi|}{\Delta\varphi_{\mathrm{disc}}}\right\rfloor,
  \qquad
  \Delta\varphi_{\mathrm{disc}} = 41.
\]
Алгебраические вехи (закрыто):

\begin{table}[htbp]
  \centering
  \caption{Вехи внутреннего спектра на $\mathbb{Z}_{512}$.}
  \label{tab:floor0-landmarks}
  \begin{tabular}{@{}rlcl@{}}
    \toprule
    метка & тики $\Phi$ & $n_E$ & смысл \\
    \midrule
    \texttt{E0\_1} & $41$ & $1$ & минимальный разрешённый скачок ($\Delta\varphi_{\min}$) \\
    \texttt{E0\_2} & $82$ & $2$ & второй уровень \\
    \texttt{pauli\_pi} & $256$ & $6$ & обменный kick $N_{\mathrm{ring}}/2$ (поворот на $\pi$) \\
    \texttt{ring\_2pi} & $512$ & $12$ & полный оборот $U(1)$ на кольце \\
    \bottomrule
  \end{tabular}
\end{table}

Связь с запретом Паули: kick на половине кольца ($256$ тиков) --- та же дискретная
$\pi$-ротация, что в конгруэнтной лестнице CL-2 (\cref{ch:evolution}).
Полный оборот $512$ тиков даёт $n_E=12$ --- двенадцать квантов $E_0$ на один цикл фазы.

\begin{proposition}[Лестница $n_E$]
\label{prop:floor0-ne-ladder}
Для $N_{\mathrm{ring}}=512$ и $\Delta\varphi_{\mathrm{disc}}=41$:
\[
  n_E(\Phi) = \lfloor |\Phi| / 41\rfloor,\quad
  \Phi_{\pi} = N_{\mathrm{ring}}/2 = 256,\quad
  \Phi_{2\pi} = N_{\mathrm{ring}} = 512.
\]
\end{proposition}

\section{Основное состояние}
\label{sec:floor0-ground}

Устойчивый планкон с $n=\pm1$ после релаксации поля (конфигурация vortex на одном~$\hv$)
имеет следующие свойства на ядре и контуре:

\begin{itemize}[noitemsep]
\item $b=1$ --- ячейка занята планконом, не планковской дыркой;
\item на контуре $|n|_{\mathrm{auto}}\ge \tfrac34$ --- устойчивый вихревой паттерн;
\item $|\Delta\varphi|_{\mathrm{core}} < \Delta\varphi_{\min}$ --- фазовый дефект сидит
  в дне Гейзенберга, не «острый'' Arg;
\item $n_E=0$ в реестре энергии: за время релаксации не накопилось
  $\ge 41$ тиков насыщающей фазы на ядре;
\item вектор Блоха на ядре устойчив под повторным применением~$g$
  (внутренняя ось $SU(2)$ не дрейфует).
\end{itemize}

\begin{proposition}[Основное состояние планкона]
\label{prop:floor0-vortex-ground}
После релаксации vortex с $n=\pm1$ на одном~$\hv$:
$b=1$, $|\Delta\varphi|_{\mathrm{core}}<\Delta\varphi_{\min}$, $n_E=0$,
устойчивый вектор Блоха на ядре.
\end{proposition}

Это \emph{нулевой} внутренний уровень по $n_E$ при ненулевом топологическом $n$:
планкон уже локализован и массивен (\cref{ch:matter}), но фазовый реестр на кольце
в основном состоянии пуст.

\section{Монодромия \texorpdfstring{$SU(2)$}{SU(2)} на ядре}
\label{sec:floor0-su2}

Топологический $n$ отвечает за обход контура; спин --- за поведение спинора \emph{на ядре}.
Вращение $R_\alpha$ вокруг оси Блоха ядра даёт скалярное произведение
\[
  \langle z,\, R_{2\pi} z\rangle \approx -1,
  \qquad
  \langle z,\, R_{4\pi} z\rangle \approx +1.
\]
Первое --- знак минус при одном обороте (фермионная монодромия);
второе --- возврат после двух оборотов.
Это тот же смысл, что у спина~$\tfrac12$ в квантовой механике,
но здесь проверяется на дискретном спиноре~$z$ в центре~$\hv$, а не постулируется.

\begin{proposition}[Полуцелый спин на ядре]
\label{prop:floor0-su2-monodromy}
Для устойчивого ядра vortex: $2\pi\mapsto -1$, $4\pi\mapsto +1$ в смысле
$\mathrm{Re}\langle z, R_\phi z\rangle$.
\end{proposition}

\section{Что ещё не закрыто}
\label{sec:floor0-open}

Следующие пункты явно открыты; они не отменяют алгебру вех и основное состояние.

\begin{table}[htbp]
  \centering
  \caption{Открытые под-leaf'ы этажа~0.}
  \label{tab:floor0-open}
  \begin{tabular}{@{}p{0.38\linewidth}p{0.52\linewidth}@{}}
    \toprule
    вопрос & статус \\
    \midrule
    возбуждения $n_E\ge 1$ в свободной эволюции &
    стандартные начальные условия не дают $|\Phi|\ge 41$ на ядре;
    нужен отдельный kick-harness \\
    полный каталог $(n_E,\,\text{Блох},\,\Phi\bmod N_{\mathrm{ring}})$ &
    конечен ($B_{hV}$), но не перечислен \\
    правила отбора между внутренними уровнями &
    схема в \cref{ch:matter}, симуляция нет \\
    внешние моды на $\varepsilon$-оболочке (1-я сфера соседей) &
    этаж~1; \S5.0.5 в модели, отдельная глава позже \\
    \bottomrule
  \end{tabular}
\end{table}

\paragraph{Связь с электроном.}
Лабораторный электрон --- частица на~$T$; планкон --- носитель на~$\Mlayer$.
Масса, заряд и радиус (через облако $\rho_\Theta$, \cref{ch:matter}) выводятся на других этажах.
Внутренний спектр этажа~0 --- \emph{зародыш} запрещёнки, общий для любого планкона,
в том числе для того, что на~$T$ читается как электрон.

\paragraph{Связь с кварками и квантовыми числами.}
Цвет, поколение, изоспин --- не здесь.
Их разбор --- после закрытия этажа~0 и внешних оболочек (этаж~1),
когда составной узел на~$\varepsilon$ даёт больше степеней свободы, чем одно~$\hv$.
